\section{Discussion}

RACK-KV is best understood as a response to the fact that long-context inference poses more than one problem at a time. If the only objective were storage, then straightforward quantization would already be compelling. If the only objective were to reduce attention work, then aggressive eviction or retrieval could be attractive. The difficulty is that long-context systems often need compressed storage, selective access, and some principled notion of safety simultaneously. The central contribution of RACK-KV is therefore not any single codec component or theorem in isolation, but the way those pieces are made compatible: historical blocks are compressed in a form that remains independently decodable, and the same blocks carry enough metadata to support query-dependent screening.

This perspective helps explain the baseline comparisons. Uniform INT8 keeps the whole cache structure intact and pays little metadata overhead, so it is a strong control for asking whether RACK-KV's additional machinery is worthwhile. The representative-layer results suggest that it is: RACK-KV spends extra bytes on anchors, indices, and block summaries, but it converts those bytes into noticeably lower attention-output error at a similar storage scale. By contrast, the matched SnapKV-style and Quest-style baselines achieve comparable budgets partly by discarding some historical information. Their higher local error is therefore not surprising. The difference is conceptual as much as numerical. RACK-KV compresses history and then skips cautiously; the selection baselines remove history directly.

The certificate's conservativeness is also easier to interpret once the storage format is part of the story. A loose theorem is not merely a mathematical inconvenience; it directly reduces the number of blocks that can be omitted after reconstruction. The observed skip rate is low because the current bound combines several worst-case relaxations: a spherical summary, exponential softmax amplification, conservative lower bounds on kept mass, and GQA sharing that prevents many logical skips from becoming physical omissions. These are not hidden failure modes. They are explicit tradeoffs in a first theorem whose primary purpose is to avoid unsafe skipping rather than to maximize skip count.

Another scientific advantage of the framework is the decomposition between compression and skipping. A heuristic retrieval score usually says only that a token or block appears important. RACK-KV instead separates the approximation into two terms: compression error from reconstructing the historical cache, and skipping error from omitting some reconstructed contributions. The theorem controls only the second term, but the decomposition clarifies exactly what remains unproved. That structure remains useful even after the initial full-model experiment: the representative-layer results explain where local distortion enters, while the teacher-forced two-prompt benchmark shows that the resulting downstream quality impact can still be small on a narrow end-to-end test.

The most important unresolved questions now shift from \emph{whether} end-to-end quality can be measured to \emph{how broadly} the current evidence should be trusted. The verified teacher-forced run covers only two short prompt categories, so it cannot settle questions about longer contexts, broader language behavior, or generation dynamics. It also leaves the systems question open. Logical certification is already enough to study safety and quality propagation, but it is not yet enough to prove lower bandwidth or faster decoding. RACK-KV should therefore be viewed as a storage and certification prototype with encouraging local and narrow end-to-end behavior rather than as a finished acceleration system.
